\documentclass[aps,prd,reprint,superscriptaddress,nofootinbib]{revtex4-2}
\usepackage{amsmath,amssymb}
\usepackage{graphicx}
\usepackage{xcolor}
\usepackage[colorlinks=true,allcolors=blue]{hyperref}

% Numbers. values/results_macros.tex carries the design values and the final numbers of the
% host-normalization and Monte Carlo accuracy analyses (black), and it defines \ph (red) with
% its colour, so \ph is not defined here. values_independent.tex carries the expected results
% of this paper (red). A sentence that ends with "% EXPECTED" changes if its number does.
\input{values/results_macros.tex}
%% values_independent.tex -- numbers of the short paper that values/results_macros.tex does not carry.
%% Input AFTER values/results_macros.tex. Every macro is \providecommand, so a value file that
%% defines the same name earlier wins. Red (\ph): expected value, to be replaced when the
%% experiment of EXPERIMENT.md finishes. Black: design value of the simulation.
%% Meaning, expected value and its source for every macro: notes/MACROS_NEEDED.md.
%% Half-widths are half the width of the 90% credible interval of H0, km/s/Mpc, 1000 mergers.

%% ---- design values not in values/MACROS.md (black) ----
\providecommand{\IndNreal}{five}         % realizations of the experiment (a word: use in text mode)
\providecommand{\IndfList}{0,\ 0.1,\ 0.3,\ 0.5\ {\rm and}\ 1}  % the input fractions of Fig. 1 (math mode)
\providecommand{\IndDmTwice}{10}         % mass-peak difference of the stronger case [Msun]
\providecommand{\IndDchiTwice}{0.20}     % mean effective spin difference of the stronger case

%% ---- Fig. 1(a): sharp mass scale, populations held at their inputs (red) ----
\providecommand{\IndHWSpec}{\ph{3.0}}          % mass spectrum alone, no catalog
\providecommand{\IndHWGal}{\ph{3.5}}           % all mergers in galaxies, galaxy catalog
\providecommand{\IndHWAGN}{\ph{0.8}}           % all mergers in AGN, AGN catalog
\providecommand{\IndHWMix}{\ph{1.7}}           % f = 0.30, labels unknown (one population, both catalogs)
\providecommand{\IndHWMark}{\ph{1.5}}          % f = 0.30, a population for each tracer, (5 Msun, 0.10)
\providecommand{\IndHWMarkTwice}{\ph{1.4}}     % f = 0.30, a population for each tracer, (10 Msun, 0.20)
\providecommand{\IndHWKnown}{\ph{1.3}}         % f = 0.30, hosts known
\providecommand{\IndHWKnownAGNPart}{\ph{1.5}}  % f = 0.30, hosts known, the mergers in AGN alone
\providecommand{\IndHWMixTenth}{\ph{2.4}}      % f = 0.10, labels unknown
\providecommand{\IndHWKnownTenth}{\ph{2.0}}    % f = 0.10, hosts known
\providecommand{\IndCostMix}{\ph{1.7}}         % (IndHWMix/IndHWKnown)^2: mergers needed, labels unknown over hosts known
\providecommand{\IndCostMark}{\ph{1.3}}        % (IndHWMark/IndHWKnown)^2
\providecommand{\IndCostMarkTwice}{\ph{1.2}}   % (IndHWMarkTwice/IndHWKnown)^2
\providecommand{\IndfAGNMix}{\ph{0.28^{+0.08}_{-0.08}}}  % f_AGN, f = 0.30, labels unknown, one realization

%% ---- Fig. 1(b): observed-like mass spectrum (red) ----
\providecommand{\IndObsHWSpec}{\ph{10}}        % mass spectrum alone
\providecommand{\IndObsHWMix}{\ph{1.9}}        % f = 0.30, labels unknown
\providecommand{\IndObsHWAGN}{\ph{0.8}}        % all mergers in AGN, AGN catalog

%% ---- Fig. 2: probability of an AGN host, one realization, f = 0.30, (5 Msun, 0.10) (red) ----
\providecommand{\IndProbAGNPos}{\ph{0.36}}     % mean over mergers truly in AGN, positions alone
\providecommand{\IndProbGalPos}{\ph{0.27}}     % mean over mergers truly in galaxies, positions alone
\providecommand{\IndProbAGNMark}{\ph{0.41}}    % mean over mergers truly in AGN, positions with masses and spins
\providecommand{\IndProbGalMark}{\ph{0.25}}    % mean over mergers truly in galaxies, positions with masses and spins

%% ---- Fig. 3: populations inferred, f = 0.30 (red) ----
\providecommand{\IndBias}{\ph{3}}              % one population fitted, (5 Msun, 0.10): combined H0 median minus input
\providecommand{\IndBiasTwice}{\ph{6}}         % the same for (10 Msun, 0.20)
\providecommand{\IndHWOnePop}{\ph{2.6}}        % one population fitted, half-width from one catalog
\providecommand{\IndOnePopCL}{\ph{99.9}}       % credible level (percent) beyond which the input H0 lies, one population, combined
\providecommand{\IndTwoPopHzero}{\ph{67.4^{+2.1}_{-2.1}}}   % a population for each tracer, one catalog: H0
\providecommand{\IndTwoPopHW}{\ph{2.1}}                     % its half-width (open symbol of Fig. 1)
\providecommand{\IndTwoPopfAGN}{\ph{0.30^{+0.06}_{-0.06}}}  % f_AGN
\providecommand{\IndTwoPopDm}{\ph{5.1^{+1.0}_{-1.0}}}       % mass-peak difference [Msun]
\providecommand{\IndTwoPopDchi}{\ph{0.10^{+0.03}_{-0.03}}}  % mean effective spin difference
\providecommand{\IndTwoPopCL}{\ph{30}}         % credible level (percent) of the input H0, a population for each tracer, combined
\providecommand{\IndLB}{\ph{9}}                % log10 Bayes factor, a population for each tracer over one population, one catalog

%% ---- Appendix C (red) ----
\providecommand{\IndRecHzeroComb}{\ph{67.9^{+0.6}_{-0.6}}}  % equal populations, f = 0.30, labels unknown, combined: H0
\providecommand{\IndRecHzeroCL}{\ph{40}}                    % credible level (percent) of the input H0 in it
\providecommand{\IndRecfAGNComb}{\ph{0.30^{+0.03}_{-0.03}}} % the same combination: f_AGN
\providecommand{\IndRecfAGNCL}{\ph{35}}                     % credible level (percent) of the input f_AGN in it
\providecommand{\IndZerofAGNUL}{\ph{0.05}}     % no merger in AGN: 90% upper limit on f_AGN, combined


\newcommand{\kms}{\,{\rm km\,s^{-1}\,Mpc^{-1}}}
\newcommand{\Msun}{\,M_\odot}
\newcommand{\fagn}{f_{\rm AGN}}
\newcommand{\dm}{\Delta m}
\newcommand{\dchi}{\Delta\chi}
\newcommand{\chieff}{\chi_{\rm eff}}

% A figure is included if its file exists; otherwise a frame of the same size holds its place.
% #1 file name without extension, #2 width, #3 height of the frame, #4 printable name
\newcommand{\showfig}[4]{%
  \IfFileExists{figures/#1.pdf}{\includegraphics[width=#2]{figures/#1.pdf}}{%
  \IfFileExists{figures/#1.png}{\includegraphics[width=#2]{figures/#1.png}}{%
  \fbox{\parbox[c][#3][c]{0.96\linewidth}{\centering\small\ttfamily figures/#4}}}}}

\begin{document}

\title{A dark siren is as good as its hosts are rare: black hole mass and spin as labels for
mergers in active galactic nuclei}

\author{Ignacio Maga\~na Hernandez}
\affiliation{McWilliams Center for Cosmology, Department of Physics, Carnegie Mellon University, Pittsburgh, Pennsylvania 15213, USA}

\date{\today}

\begin{abstract}
Thousands of galaxies fill every localization volume, so dark sirens take most of their Hubble
constant $H_0$ from the mass spectrum. Active galactic nuclei (AGN) are a hundred times rarer.
In simulations of $\MockNev$ binary black hole mergers with a known, sharp mass scale, the 90\%
half-width of $H_0$ is $\IndHWGal\kms$ with all mergers in galaxies, $\IndHWAGN$ with all in
AGN. % EXPECTED
With $\MockfAGN$ in AGN, unidentified, a two-tracer likelihood gives $\IndHWMix$; if those are
heavier by $\MockDmuG\Msun$ and spin faster by $\MockDmuChi$, their masses and spins label them
and give $\IndHWMark$, while one fitted distribution for all mergers biases $H_0$ by
$\IndBias$. % EXPECTED
It takes a complete AGN catalog and a distribution for each tracer: a dark siren is as good as
its hosts are rare.
\end{abstract}

\maketitle

%% ------------------------------------------------------------------------------------------
%% Introduction (no heading)
%% ------------------------------------------------------------------------------------------

A binary black hole merger gives its luminosity distance but no redshift. The dark siren
method takes the redshift from the galaxies in the localization
volume~\cite{Schutz:1986gp,Gray:2019ksv,Gair:2022zsa}, and its weak point is their number: a
volume that holds thousands of galaxies returns little more than the smooth redshift distribution
of all galaxies. The Hubble constant $H_0$ comes mostly from the black hole mass spectrum, whose
features supply a redshift without any
catalog~\cite{Taylor:2011fs,Farr:2019twy,Ezquiaga:2022zkx,Mastrogiovanni:2021wsd}. The fifth
gravitational-wave transient catalog gives $H_0=71.7^{+17.7}_{-11.4}\kms$ (90\% credible
interval, with the bright siren GW170817) from 236 sources, and its galaxy catalog improves the
constraint by 3.5\% over the mass spectrum alone (Ref.~\cite{LIGOScientific:2026uyd}, version 2
of the preprint).

Fewer candidate hosts would change this. The accretion disks of active galactic nuclei (AGN)
are a proposed site of black hole
mergers~\cite{Bartos:2016dgn,Stone:2016wzz,Mckernan:2017ssq}, and luminous AGN are far rarer
than galaxies. In the simulated catalogs of this paper, with $\MocknAGN$ AGN and $\MocknGal$
galaxies per ${\rm Mpc^{3}}$, the median localization volume holds $\PureAGNHostsPerVol$ AGN and
$\PureGalHostsPerVol$ galaxies (Fig.~\ref{fig:concept}a). A merger in an AGN, matched to an AGN
catalog, therefore sits between a dark siren and a bright one. $H_0$ has been inferred together
with the fraction of mergers in
AGN~\cite{Bom:2023zgw,Cross-Parkin:2026pua}, and that fraction has been inferred at fixed
cosmology~\cite{Veronesi:2023ugk,Veronesi:2024tef,Zhu:2025rdv,Pouw:2026lar}. None of these
analyses infers a mass or spin distribution that depends on the host, although black holes that
merge in AGN disks are expected to be heavier and to spin
faster~\cite{Yang:2019cbr,Tagawa:2020dxe,Tagawa:2021ofj}; Refs.~\cite{Bom:2023zgw,Pouw:2026lar,Cross-Parkin:2026pua}
leave that inference to future work.

\begin{figure*}
\showfig{fig_concept}{\textwidth}{0.34\textwidth}{fig\_concept}
\caption{The two tracers and the two populations in one realization of the simulation. (a) A
slice of the simulated catalogs within $\IndConceptDec^\circ$ of the equator and below
$z=\IndConceptzMax$, in comoving distance (bar: $\IndConceptScale$~Mpc): $\IndConceptNGal$
objects of the galaxy catalog (blue dots) and $\IndConceptNAGN$ of the AGN catalog (vermilion
diamonds). The black outline is the 90\% localization region of one AGN-hosted merger at
$z=\IndConceptzTrue$, seen from the side in the zoom; it holds $\IndConceptVolNGal$ galaxies and
$\IndConceptVolNAGN$ AGN, where the median region holds $\PureGalHostsPerVol$ and
$\PureAGNHostsPerVol$, and the ring marks the true host of the merger. (b) Source-frame
primary mass and (c) effective spin of the $\MockNev$ detected mergers of a realization in which
the populations differ, $\IndConceptNgalHosted$ galaxy-hosted and $\IndConceptNAGNHosted$
AGN-hosted (true values, each class normalized to unit area); dashed lines mark the input mass
peaks, $\MockMuG$ and $\MockMuGAGN\Msun$, and the input spin means, $\MockMuChi$ and
$\MockDmuChi$. An AGN-hosted merger has a handful of possible hosts where a galaxy-hosted one
has hundreds, and its mass and spin, which overlap with those of the others, say in which
catalog to look.}
\label{fig:concept}
\end{figure*}

We ask where the $H_0$ information of a catalog of mergers sits when some of them are in AGN,
and what the masses and spins of the black holes add to it. The rate of mergers with source
masses and effective spin $\theta$, at redshift $z$ and sky position $\Omega$, is a sum over the
two tracers of large-scale structure,
\begin{equation}
\frac{dN}{d\theta\,dz\,d\Omega}\;\propto\sum_{k\,=\,{\rm gal},\,{\rm AGN}}
f_k\;p_k(z,\Omega)\;p_{\rm pop}(\theta\,|\,\Lambda_k),
\label{eq:rate}
\end{equation}
where $\fagn$ is the fraction of mergers hosted by AGN, $f_{\rm gal}=1-\fagn$, $p_k$ is the
distribution of the hosts of tracer $k$, built from its catalog and normalized to unity, and
$\Lambda_k$ are the parameters of the mass and spin distribution (the population) of the mergers
in tracer $k$. This is the merger rate per galaxy of Ref.~\cite{Mastrogiovanni:2023emh} written
for each tracer; that work notes that the population may depend on the host and sets the
dependence aside. Appendix~\ref{app:like} gives the likelihood with its selection term. Every
result below comes from simulations (Appendix~\ref{app:sim}): galaxies and AGN drawn from one
density field, with redshift errors $\MockDz(1+z)$, and $\MockNev$ detected mergers per
realization, with input $H_0=\MockHzero\kms$ in flat $\Lambda$CDM. Unless stated, the catalogs
are complete and the
populations are held at their inputs. Priors are uniform, intervals are 90\% credible, and the
half-width is half the width of the 90\% interval of $H_0$ from $\MockNev$ mergers.

%% ------------------------------------------------------------------------------------------
\section{How much is a host catalog worth?}
\label{sec:worth}
%% ------------------------------------------------------------------------------------------

With every merger in an AGN and matched to the AGN catalog ($\fagn=1$ in Eq.~\eqref{eq:rate}),
the half-width is $\IndHWAGN\kms$; with every merger in a galaxy and matched to the galaxy
catalog ($\fagn=0$), it is $\IndHWGal$. % EXPECTED
The galaxy catalog is worth nothing here: the mass spectrum alone, with both $p_k$ replaced by a
uniform distribution in comoving volume, gives $\IndHWSpec$ (Fig.~\ref{fig:budget}a). % EXPECTED

The number of candidate hosts decides (Fig.~\ref{fig:concept}a). The $\PureGalHostsPerVol$
galaxies of a localization
volume overlap many times within their redshift errors, and their sum is close to the volume
element that the mass spectrum analysis already uses. The $\PureAGNHostsPerVol$ AGN are a sparse
set of redshifts, different for each merger, and only near the true $H_0$ do a thousand
distances all find a host.

Each number carries a choice of the simulation. The $\IndHWSpec$ belongs to a mass spectrum with
$\MockPeakFrac$ of the mergers in a Gaussian peak at $\MockMuG\Msun$, chosen to make the mass
scale as informative as it can be; for a spectrum like the fit to the third transient
catalog~\cite{KAGRA:2021duu}, with $\ObsPeakFrac$ of the mergers in the peak, the mass spectrum
alone gives $\IndObsHWSpec$ (Fig.~\ref{fig:budget}b). % EXPECTED
The $\IndHWAGN$ belongs to an AGN density of $\MocknAGN\,{\rm Mpc^{-3}}$ and to redshift errors
of $\MockDz(1+z)$. It does not rest on the mass scale: the observed-like spectrum gives
$\IndObsHWAGN$. % EXPECTED
The Monte Carlo sums weigh on it: redrawing the posterior samples moves the median of $H_0$ by
$\MCAGNHzeroSdEightK\kms$ (Appendix~\ref{app:checks}).

\begin{figure*}
\showfig{fig1_budget}{\textwidth}{0.34\textwidth}{fig1\_budget}
\caption{The 90\% half-width of the $H_0$ posterior from $\MockNev$ simulated mergers against the
fraction $\fagn$ of mergers hosted by AGN, with complete catalogs and the populations held at
their inputs (means over \IndNreal{} realizations). Labels unknown: both catalogs, $\fagn$ free,
one population for all mergers. A population for each tracer: the same, for mergers in AGN that
are heavier by $\dm=\MockDmuG\Msun$ and spin faster by $\dchi=\MockDmuChi$, and for twice those
values. Hosts known: each merger matched to the catalog of its true tracer. Horizontal lines: the
mass spectrum alone, with no catalog, and the galaxy catalog alone when no merger is in an AGN.
The open symbol marks a population for each tracer inferred with $H_0$. (a) A mass spectrum with
$\MockPeakFrac$ of the mergers in its Gaussian peak. (b) A spectrum like the observed one, with
$\ObsPeakFrac$ in the peak. With three mergers in ten in AGN the half-width in (a) falls from
$\IndHWSpec$ to $\IndHWMix\kms$ with labels unknown, and masses and spins recover part of the
way to hosts known ($\IndHWKnown$).} % EXPECTED
\label{fig:budget}
\end{figure*}

%% ------------------------------------------------------------------------------------------
\section{Which mergers are in AGN?}
\label{sec:which}
%% ------------------------------------------------------------------------------------------

The likelihood finds them statistically, at a cost. With $\MockfAGN$ of the mergers in AGN and
no label that says which (labels unknown), Eq.~\eqref{eq:rate} with $\fagn$ free and one
population for all mergers gives a half-width of $\IndHWMix\kms$, against $\IndHWKnown$ when the
tracer of every merger is given (hosts known); to match the latter it needs $\IndCostMix$ times
as many mergers (Fig.~\ref{fig:budget}a). % EXPECTED

The second number shows where the information sits. With hosts known, the mergers in AGN alone
give $\IndHWKnownAGNPart\kms$, and the other seven in ten, with their catalog and their mass
spectrum, bring it only to $\IndHWKnown$. % EXPECTED
Three mergers in ten carry most of the measurement. % EXPECTED

With labels unknown the likelihood has to find those mergers by their positions. The mergers in
AGN fall on AGN more often than chance allows, which measures $\fagn=\IndfAGNMix$ for an input
of $\MockfAGN$. But the two tracers follow the same
structure: the localization volume of a merger in a galaxy holds $\PureAGNHostsPerVolGalEvents$
AGN, against $\PureAGNHostsPerVol$ for a merger in an AGN. From positions alone the mergers
truly in AGN receive a mean probability of $\IndProbAGNPos$ of being in an AGN and those truly
in galaxies $\IndProbGalPos$, where $\MockfAGN$ for both would mean no information
(Fig.~\ref{fig:labels}). % EXPECTED

A small fraction of mergers in AGN is enough to matter. % EXPECTED
With $\MockfAGNTenth$ of them in AGN the half-width is $\IndHWMixTenth\kms$ with labels unknown
and $\IndHWKnownTenth$ with hosts known, against $\IndHWGal$ when no merger is in an
AGN. % EXPECTED
For the spectrum like the observed one, $\MockfAGN$ in AGN with labels unknown gives
$\IndObsHWMix$, against $\IndObsHWSpec$ from the mass spectrum alone: there the AGN catalog
supplies nearly all of the measurement (Fig.~\ref{fig:budget}b). % EXPECTED

\begin{figure}
\showfig{fig2_labels}{\columnwidth}{0.62\columnwidth}{fig2\_labels}
\caption{Probability that each merger is hosted by an AGN, for the mergers truly in galaxies and
those truly in AGN of one simulated catalog ($\MockNev$ mergers, $\fagn=\MockfAGN$,
$\dm=\MockDmuG\Msun$, $\dchi=\MockDmuChi$). Positions alone: one population for all mergers.
Positions with masses and spins: a population for each tracer. The dashed line is $\fagn$, the
value for a merger that says nothing about its host. Masses and spins raise the mean probability
of the mergers truly in AGN from $\IndProbAGNPos$ to $\IndProbAGNMark$, yet no single merger is
identified: the labels act on the catalog as a whole.} % EXPECTED
\label{fig:labels}
\end{figure}

%% ------------------------------------------------------------------------------------------
\section{Can the black holes tell us?}
\label{sec:labels}
%% ------------------------------------------------------------------------------------------

In part. % EXPECTED
If mergers in AGN have the peak of the mass spectrum higher by $\dm=\MockDmuG\Msun$ and the mean
effective spin higher by $\dchi=\MockDmuChi$ (Fig.~\ref{fig:concept}b,~c), a population for each
tracer in
Eq.~\eqref{eq:rate} gives a half-width of $\IndHWMark\kms$, against $\IndHWMix$ when the
populations are equal; for twice those differences it gives $\IndHWMarkTwice$, against
$\IndHWKnown$ with hosts known (Fig.~\ref{fig:budget}a). % EXPECTED

Masses and spins act as labels. With them the mean probability of being in an AGN rises from
$\IndProbAGNPos$ to $\IndProbAGNMark$ for the mergers truly in AGN and falls from
$\IndProbGalPos$ to $\IndProbGalMark$ for those truly in galaxies
(Fig.~\ref{fig:labels}). % EXPECTED
No single merger is identified: the two populations overlap (Fig.~\ref{fig:concept}b,~c), and
at the detection threshold the effective spin is measured to $\MockSigChi$, twice the difference
between them. % EXPECTED
The labels act on the catalog as a whole, and the cost of unknown hosts falls from $\IndCostMix$
times as many mergers to $\IndCostMark$, and to $\IndCostMarkTwice$ for twice the
differences. % EXPECTED

These numbers hold both populations at their inputs. With the peak and the spin distribution of
each tracer inferred along with $H_0$ and $\fagn$, the half-width at $\fagn=\MockfAGN$ is
$\IndTwoPopHW\kms$ (open symbol in Fig.~\ref{fig:budget}a), against $\IndHWMark$; we read the
loss as the location of the peak and $H_0$ both setting the redshifted mass scale. % EXPECTED

How different the two populations are is open. The analyses that use AGN catalogs stop short of
inferring the difference: Ref.~\cite{Pouw:2026lar} gives both channels the same mass and
spin distributions, Ref.~\cite{Cross-Parkin:2026pua} expects host sub-populations to differ in
mass and spin, and Ref.~\cite{Bom:2023zgw} simulates mergers above $35\Msun$ in AGN with a
likelihood that contains no mass distribution.

%% ------------------------------------------------------------------------------------------
\section{Is one population for all mergers enough?}
\label{sec:bias}
%% ------------------------------------------------------------------------------------------

Not if the populations differ. One mass and spin distribution, with its peak and spin parameters
fitted along with $H_0$ and $\fagn$, puts $H_0$ high by $\IndBias\kms$ for $\dm=\MockDmuG\Msun$ and
$\dchi=\MockDmuChi$, and by $\IndBiasTwice$ for twice those values, against a half-width of
$\IndHWOnePop$ from one catalog of $\MockNev$ mergers (Fig.~\ref{fig:bias}a). % EXPECTED

Fitting the single population does not protect $H_0$. We read the bias as an effect of
selection: heavier mergers are heard farther away, so the mergers in AGN are $\MockfAGN$ of all
mergers and $\IndConceptNAGNHosted$ of the $\MockNev$ detected ones in
Fig.~\ref{fig:concept}, and their share grows with distance. One
population can account for masses that rise with distance only by more redshift per unit
distance, a higher $H_0$. % EXPECTED
The size of the bias belongs to a mass spectrum with $\MockPeakFrac$ of the mergers in the peak,
where the mass scale weighs most.

A population for each tracer removes the bias and measures the difference. With $\dm$ and
$\dchi$ free, one catalog gives $H_0=\IndTwoPopHzero\kms$, $\fagn=\IndTwoPopfAGN$,
$\dm=\IndTwoPopDm\Msun$ and $\dchi=\IndTwoPopDchi$, for inputs of $\MockHzero$, $\MockfAGN$,
$\MockDmuG$ and $\MockDmuChi$ (Fig.~\ref{fig:bias}b). % EXPECTED
Combined over \NrealHD{} realizations, the input $H_0$ lies at the $\IndTwoPopCL\%$ credible
level of this analysis; combined over \NrealHDSingle{}, it lies beyond the $\IndOnePopCL\%$
level of the one-population analysis. % EXPECTED
One catalog favors the two populations with a Bayes factor of $\log_{10}B=\IndLB$, for priors
uniform on $\PriorDmuGLo<\dm<\PriorDmuGHi\Msun$ and
$\PriorDmuChiLo<\dchi<\PriorDmuChiHi$. % EXPECTED
The tie between host and black hole is a gain when it is modeled and a bias when it is ignored,
so an analysis that brings in an AGN catalog should let the population depend on the
tracer. % EXPECTED

\begin{figure*}
\showfig{fig3_bias}{\textwidth}{0.34\textwidth}{fig3\_bias}
\caption{Mergers whose populations differ between the tracers ($\dm=\MockDmuG\Msun$,
$\dchi=\MockDmuChi$, $\fagn=\MockfAGN$), analyzed with the populations inferred. (a) Posterior
density of $H_0$ combined over realizations of $\MockNev$ mergers, with one population fitted to
all mergers (\NrealHDSingle{} realizations) and with a population for each tracer (\NrealHD{});
the vertical line is the input. (b) Joint posterior of the two population differences in the
second analysis, with its 90\% credible contour; the cross is the input and the circle is zero,
equal populations. One population puts $H_0$ high by $\IndBias\kms$; a population for each
tracer removes the bias and measures both differences.} % EXPECTED
\label{fig:bias}
\end{figure*}

%% ------------------------------------------------------------------------------------------
\section{What does it take?}
\label{sec:take}
%% ------------------------------------------------------------------------------------------

For $H_0$, an AGN catalog complete over the volume of the detected mergers, here
$z<\MockzHorizon$, and nothing else: with complete catalogs no abundance of AGN enters
Eq.~\eqref{eq:rate}. For $\fagn$, also the abundance of AGN: at the shallowest flux limit we
consider, where the catalogs hold the hosts of $\MockEventCompEighteen\%$ of the mergers, an AGN
density given to the model at twice its true value moves $\fagn$ from $\NormfEighteenProd$ to
$\NormfEighteenProdTwo$ and $H_0$ from $\NormHzeroEighteenProd$ to
$\NormHzeroEighteenProdTwo\kms$.

The reason is that mergers measure a rate per host. Each $p_k$ is normalized by the number of
hosts $N_k$, so what the data constrain is the merger rate per AGN relative to the rate per
galaxy, $[\fagn/(1-\fagn)]\,N_{\rm gal}/N_{\rm AGN}$, which we find to be
$\NormRateRatioCompleteProd$ for an input of $\MockRateRatio$. The fraction follows once the
hosts are counted. A catalog that misses hosts needs their number density, and it supplies that
itself where it is complete, here to $\MocknAGNScatterEighteen\%$, which moves $\fagn$ by
$\NormDeltafCountsScatterEighteen$. The association searches meet the same limit as a product of
$\fagn$ and the completeness of the AGN
catalog~\cite{Veronesi:2023ugk,Veronesi:2024tef,Zhu:2025rdv}, and Ref.~\cite{Pouw:2026lar} finds
from simulated catalogs that the catalog must be complete within the localization volumes.
Flux-limited catalogs, and number densities inferred with the mergers, are the subject of a
companion paper.

Two quantities decide how much of Fig.~\ref{fig:budget} applies to the real catalogs. The
fraction of mergers in AGN has an upper limit of 0.21 (95\%) for AGN with bolometric luminosity
above $10^{44.5}\,{\rm erg\,s^{-1}}$~\cite{Veronesi:2024tef}, and one analysis reports
$0.39^{+0.41}_{-0.32}$ (90\%)~\cite{Zhu:2025rdv}; at $\MockfAGNTenth$ the half-width is already
$\IndHWMixTenth\kms$, against $\IndHWGal$ with galaxies alone. % EXPECTED
The AGN catalog has to be complete to the distance of the detected mergers, a requirement on
all-sky AGN samples such as Quaia~\cite{Storey-Fisher:2023gca}. Given both, the analysis in use
needs two additions, $\fagn$ and a population for each tracer, and the black holes will say
which hosts are theirs. A dark siren is as good as its hosts are rare.

%% ==========================================================================================
\appendix
%% ==========================================================================================

\section{The likelihood}
\label{app:like}

Merger $i$ has $n=\ProdNsamp$ posterior samples $\theta_{ij}$ of the detector-frame primary mass
$m_1^{\rm d}$, the mass ratio $q$, the luminosity distance $d_L$, the effective spin $\chieff$
and the sky position $\Omega$, drawn under the prior $\pi_{\rm PE}$. For parameters
$\Lambda=(H_0,\fagn,\Lambda_{\rm gal},\Lambda_{\rm AGN})$, Eq.~\eqref{eq:rate} in these
coordinates is
\begin{equation}
\pi(\theta\,|\,\Lambda)=\sum_k f_k\,
\frac{p_{\rm pop}(m_1,q,\chieff\,|\,\Lambda_k)\;p_k(z,\Omega)}{(1+z)^2\,dd_L/dz},
\label{eq:weight}
\end{equation}
with $z=z(d_L;H_0)$ and $m_1=m_1^{\rm d}/(1+z)$. One factor $1+z$ is the time dilation of a
merger rate per host that is constant in source-frame time; the other, with $dd_L/dz$, is the
Jacobian from $(m_1,z)$ to $(m_1^{\rm d},d_L)$. The hosts of tracer $k$ in the sky pixel of
$\Omega$ (HEALPix, $N_{\rm side}=\MockNside$) are distributed as
\begin{equation}
p_k(z,\Omega)=\frac{1}{N_k}\sum_{g\,\in\,{\rm pix}(\Omega)}
\frac{\mathcal{N}(z;z_g,\sigma_g)\,dV_c/dz}
{\int\mathcal{N}(z';z_g,\sigma_g)\,(dV_c/dz')\,dz'},
\label{eq:hosts}
\end{equation}
a sum over the cataloged objects $g$ of the pixel with redshifts $z_g$ and errors
$\sigma_g=\MockDz(1+z_g)$, each with the same weight, so that every host of a tracer has the
same merger rate; $N_k$ is the number of objects in the catalog. A flux-limited catalog adds to
the sum and to $N_k$ the hosts it misses out to $z=\NormZmaxProd$, beyond every detected merger:
$n_k[1-C_k(z)]\,dV_c$, with $n_k$ the number density and $C_k$ the completeness. This is how an
abundance enters (Sec.~\ref{sec:take}).

With the merger rate marginalized over a prior uniform in its
logarithm~\cite{Loredo:2004nn,Mandel:2018mve}, the likelihood of $N$ detected mergers is
\begin{equation}
\ln\mathcal{L}=\sum_{i=1}^{N}\Big[\ln\hat Z_i+\frac{\sigma_i^2}{2}\Big]-N\ln\hat\mu
+\frac{N(N+3)}{2N_{\rm eff}},
\label{eq:like}
\end{equation}
\begin{equation}
\hat Z_i=\frac1n\sum_{j=1}^{n}\frac{\pi(\theta_{ij}|\Lambda)}{\pi_{\rm PE}(\theta_{ij})},\qquad
\hat\mu=\frac{1}{N_{\rm draw}}\sum_{l\,\in\,{\rm det}}\frac{\pi(\theta_l|\Lambda)}{p_{\rm draw}(\theta_l)},
\label{eq:sums}
\end{equation}
where $\hat\mu$ is the detected fraction, summed over the detected ones of $N_{\rm draw}$
simulated signals drawn from $p_{\rm draw}$. Both sums are Monte Carlo averages. The logarithm
of an average is low by half its relative variance, which for $\hat Z_i$ is
$\sigma_i^2=\sum_j w_{ij}^2/(\sum_j w_{ij})^2-1/n$ with $w_{ij}$ the terms of the sum, and the
bracket of Eq.~\eqref{eq:like} restores it. The last term marginalizes over the uncertainty of
$\hat\mu$, whose effective number of simulated signals is
$N_{\rm eff}$~\cite{Farr:2019rap}, and we require
$N_{\rm eff}>\NeffFactor N$~\cite{Farr:2019rap,Essick:2022ojx,Talbot:2023pex}.
% The sigma_i^2/2 term is written from the settings list of the experiment ("half-variance
% correction of ln L"); confirm its exact form against the code before submission.

The analyses of Fig.~\ref{fig:budget} are limits of
Eqs.~\eqref{eq:weight}--\eqref{eq:like}. One population for all mergers sets
$\Lambda_{\rm AGN}=\Lambda_{\rm gal}$. The galaxy catalog alone sets $\fagn=0$. The mass
spectrum alone replaces both $p_k$ by $dV_c/dz$, normalized, with no dependence on $\Omega$.
With hosts known, each merger keeps only the term of its true tracer, and the likelihood is the
product of one factor for the mergers in galaxies and one for the mergers in AGN, each with its
own $\hat\mu$. The probability that merger $i$ is in an AGN (Fig.~\ref{fig:labels}) is the share
of the AGN term in $\hat Z_i$ at the posterior median of $H_0$ and $\fagn$.

The matter density is fixed at $\Omega_m=\MockOmegam$. The priors are uniform, on
$\PriorHzeroLo<H_0<\PriorHzeroHi\kms$ and $\PriorfLo<\fagn<\PriorfHi$ and, where the populations
are inferred, on the peak location ($\PriorMuGLo$ to $\PriorMuGHi\Msun$) and width
($\PriorSigmaGLo$ to $\PriorSigmaGHi\Msun$), the mean ($\PriorMuChiLo$ to $\PriorMuChiHi$) and
width ($\PriorSigmaChiLo$ to $\PriorSigmaChiHi$) of the effective spin, and the two differences
(Sec.~\ref{sec:bias}); the other population parameters stay at their inputs. Posteriors of
$(H_0,\fagn)$ are evaluated on grids, and the inferred populations are sampled with nested
sampling~\cite{Speagle:2019ivv}, which gives the evidence.

\section{The simulated catalogs and mergers}
\label{app:sim}

The catalogs are Poisson samples of one lognormal density
field~\cite{Tessore:2023zyk,Lewis:1999bs} in shells of $\MockShellWidth$~Mpc out to
$z=\MockzCat$: galaxies with bias $\MockbGal$ and comoving density
$\MocknGal\,{\rm Mpc^{-3}}$, that of galaxies brighter than $\MockLcut\,L_*$, and AGN with bias
$\MockbAGN$ and density $\MocknAGN\,{\rm Mpc^{-3}}$, that of AGN with 2 to 10~keV luminosity
above $10^{\MockLXcut}\,{\rm erg\,s^{-1}}$. They hold $\MockNGal$ galaxies and $\MockNAGN$ AGN.
Each cataloged redshift carries a Gaussian error of $\MockDz(1+z)$.

A fraction $\fagn$ of the mergers ($\IndfList$) is placed on AGN and the rest on galaxies,
uniformly over the objects of each catalog, with a rate per host constant in source-frame time.
Primary masses follow a power law of index $-\MockAlpha$ between $\MockMmin$ and
$\MockMmax\Msun$ plus a Gaussian peak~\cite{Talbot:2018cva} at $\MockMuG\Msun$ of width
$\MockSigmaG\Msun$ that holds $\MockPeakFrac$ of the mergers and $\MockDetPeakShare$ of the
detected ones; mass ratios follow a power law of index $\MockBeta$ and effective spins a
Gaussian of mean $\MockMuChi$ and width $\MockSigmaChi$. The observed-like spectrum takes its
mass parameters from the fit to the third transient catalog~\cite{KAGRA:2021duu}: index
$-\ObsAlpha$, maximum mass $\ObsMmax\Msun$, and a peak at $\ObsMuG\Msun$ of width
$\ObsSigmaG\Msun$ with $\ObsPeakFrac$ of the mergers and $\ObsDetPeakShare$ of the detected
ones. Where the populations differ, mergers in AGN have the peak at $\MockMuGAGN\Msun$ and a
mean effective spin of $\MockDmuChi$, or $\dm=\IndDmTwice\Msun$ and $\dchi=\IndDchiTwice$ in the
stronger case.

A merger is detected when its observed signal-to-noise ratio $\rho$, the optimal one plus a unit
Gaussian, exceeds $\MockSNRthr$; the optimal one is $\MockRhoRef$ for a redshifted chirp mass of
$\MockMcRef\Msun$ at 1~Gpc. The errors scale as $\MockSNRthr/\rho$: $\MockSigLnMc$ in the
logarithm of the chirp mass, $\MockSigLnq$ in that of the mass ratio and $\MockSigChi$ in the
effective spin, with a sky error between $\MockSkyMin$ and $\MockSkyMax$ degrees; the distance
follows from $\rho$ and the chirp mass, with a median error of $\MockSigLnDL\%$. The posterior
samples are exact for this measurement, and detection depends on the observed data
alone~\cite{Essick:2023upv}. Detected mergers reach $z=\MockzHorizon$, with a median of
$\MockzMedian$. Each realization has its own catalogs, its own $\MockNev$ mergers and its own
simulated signals for $\hat\mu$, $\MockNinjFound$ detected of $\MockNinjDrawn$.

\section{Recovery of the inputs and accuracy of the Monte Carlo sums}
\label{app:checks}

With equal populations, $\fagn=\MockfAGN$ and labels unknown, the posterior combined over
\NrealKnown{} realizations is $H_0=\IndRecHzeroComb\kms$ and $\fagn=\IndRecfAGNComb$, with the
inputs at its $\IndRecHzeroCL\%$ and $\IndRecfAGNCL\%$ credible levels; with no merger in AGN,
\NrealCtrl{} realizations give $\fagn<\IndZerofAGNUL$ (90\%). % EXPECTED
Redrawing the $\ProdNsamp$ posterior samples of every merger moves the median of $H_0$ by
$\MCMixedHzeroSdEightK\kms$ (standard deviation over $\MCMixedNEightK$ redraws) with $\MockfAGN$
of the mergers in AGN and by $\MCAGNHzeroSdEightK$ with all of them in AGN, against half-widths
of $\IndHWMix$ and $\IndHWAGN$; another set of simulated signals moves it by
$\MCMixedInjHzeroSd$. % EXPECTED
With $\MCNsampLo$ samples the width of the 90\% interval itself varies by $\MCWidthChange\%$
between redraws.

\bibliography{refs}

\end{document}
